\section{DISCUSSION}\label{sec:discussion}

\subsection{Implications for Evaluating Governed Agents}
The mock-versus-live divergence (Section~\ref{sec:livemock}) is itself a claim about how systems like
this one should be evaluated, and four practices follow from the results and lessons of
Section~\ref{sec:lessons}. First, state whether the model in the loop was real or mocked, and give
both numbers whenever both exist, since a mock figure describes only the mock. Second, include cheap
non-LLM baselines: two idealized ones beat the full agent system on recovery time and manual
interventions, a fact a comparison against static orchestration alone would conceal. Third, grade a
safety claim against an oracle that shares no code with the mechanism under test, since the
independent oracle here caught two contract-layer gaps the system's own tests had missed. Fourth, run
a short, paid pilot at the real settings before committing to a long campaign; ours caught a retired
model on its very first call.

\subsection{Deployment Implications}
The decision-quality figure most relevant to deployment is $\NDeltaDecisionCorrectFull{}$, not the
headline MTTR number: at these settings \texttt{full} executes an accepted mitigation in under half
its runs. Taken together with the rise in manual interventions, this campaign offers no support for
running this class of system unattended and fully autonomous against real infrastructure. It instead
supports the graduated trust ladder of Section~\ref{sec:system}, shadow, approval, and autonomous
modes plus a durable kill switch and a per-target action cap, as the safer production default, and it
argues against treating the research benchmark's default of full autonomy as a deployment
recommendation. For now, the trust core is an architectural response to a question our own
measurements raise; we have not yet measured whether it actually works.

\subsection{The Manual-Intervention Mechanism Is Measured, Not Explained}
Section~\ref{sec:results} reports that \texttt{full} and \texttt{recovery\_only} increase manual
interventions relative to \texttt{baseline}, contrary to the base paper's claim. We confirmed the
measurement itself: the metric is a direct count of rows in the
\texttt{telemetry.\allowbreak manual\_\allowbreak interventions} table that every configuration
writes to. We did not confirm the mechanism behind it. Repeated escalation across ticks of one open
fault could dominate, or rate-limit and budget denials that turn into human hand-offs, or some
interaction of the two. Telling them apart would need a per-tick trace of proposal outcomes across
the fault's lifetime, a grain the current telemetry does not retain. We leave the mechanism open
rather than infer it from the aggregate count, and a reader should not treat it as resolved in the
system's favor.

\subsection{Limitations and Threats to Validity}\label{sec:limitations}

\subsubsection{Threats to Validity}
Six threats apply directly. \emph{Simulated human.} Every comparison against \texttt{baseline}
inherits the log-normal on-call model; Section~\ref{sec:sensitivity} reports its break-even rather
than defending the model itself. \emph{Modelled baselines.} \texttt{rule\_based} and
\texttt{autoscale} resolve covered faults after fixed, assumed latencies (Section~\ref{sec:method}),
so their edge over the agents is a statement about those assumptions, not a measurement of a deployed
rule engine or autoscaler. \emph{Synthetic workload.} The data come from a downscaled synthetic table
with injected faults, so we have no real production traces, and the results may not carry over to
them. \emph{One live model family.} The live arm used a single reasoning model and a single faster
model from one provider at temperature zero, so the decision-quality and cache-poisoning findings
(Section~\ref{sec:results}) are properties of that particular pairing and may not generalize.
\emph{Unequal replicate counts.} \texttt{baseline}, \texttt{rule\_based}, \texttt{autoscale}, and
\texttt{full} each carry $\NNFull{}$ replicates, while the four single-agent ablations carry half
that, a cost trade-off made before the campaign (Section~\ref{sec:method}) that widens their
confidence intervals relative to the headline comparisons. \emph{Nominal seed matching.} Because the
per-run seed is a function of the configuration name (Section~\ref{sec:method}), no two
configurations share fault or human-latency randomness at a matched (scenario, replicate) index. We
therefore used the unpaired rank-sum test as primary, and the retained signed-rank statistic should
not be treated as evidence from a genuinely paired design.

\subsubsection{Scope of the Testbed}
The testbed is one machine. The resource-contention scenario applies real CPU stress, which is also
why runs cannot be parallelized and why the campaign's size is bounded. The policy engine is pinned
to one release, and lesson L6 (Section~\ref{sec:lessons}) shows that a newer release from the same
vendor would have changed the adversarial result had we not pinned it; we disclose this as a property
of our evaluation rather than a defect we investigated in the policy engine itself. Only four fault
types were injected, while a real pipeline's incident taxonomy runs larger, and lesson L4
(Section~\ref{sec:lessons}) shows that mapping an injected fault type onto a real incident type is
not free.

\subsubsection{Evidence Gaps}
Two contributions described in Section~\ref{sec:system} are implemented but not evaluated, and a
third gap sits inside a contribution that otherwise is evaluated. \emph{Cross-model generalization.}
The harness that could test the base paper's model-agnostic claim is unit-tested, but its live sweep
against real provider models has never been run and recorded, so nothing here counts as evidence for
or against model-agnosticism. \emph{The trust core.} Graduated autonomy, the kill switch, and the
blast-radius cap are implemented and unit-tested, yet every configuration ran under one fixed mode,
leaving no mode transition and no kill-switch-to-halt latency measured. \emph{Cost sub-components.}
The disclosed cost model of Eq.~\ref{eq:cost} computes a compute/storage split internally (Section~\ref{sec:results}), but
the campaign's committed per-run harvest keeps only the combined total, so the split cannot be
reported from this evidence even though the total can.

Each gap has a concrete, bounded remedy: run the cross-model sweep against at least two live
providers on the same scenario set; run a dedicated drill that engages the kill switch and a mode
transition under load and logs the timing; extend the metrics harvest to persist the cost ledger's
component breakdown. None of the three is a design flaw. Each is simply evidence this campaign did
not collect.

\subsubsection{What Would Change Our Conclusion}
A cross-model repetition of Section~\ref{sec:livemock} that reproduced the same order-of-magnitude
mock-versus-live gap with a different provider would strengthen the central claim considerably. One
that did not would suggest the gap is specific to this model pairing rather than a property of mock
evaluation as a method, and it would deserve just as plain a report.
